% !TeX root = ../main.tex
\section{Evaluation Protocol}
\label{sec:evaluation}

\subsection{Compared Models}
\label{subsec:compared_models}

We compare four models from the measured-mass-properties branch. The calibrated aerodynamic model is the DeLaurier-style prior after train-only affine force alignment of the exported prior predictions. The first correction model applies a force-only structured correction and leaves the prior moment channels unchanged; it tests whether real-flight force residuals can be reduced without introducing a moment model. The deployable correction extends this model with simulator-available rates, airdata, phase, and control features for both force and effective-moment prediction. The phase-structured correction uses the same deployable feature boundary and adds explicit wingbeat-phase harmonic and interaction terms, giving the final correction used for residual-structure analysis.

Together, these models separate the effects of affine prior alignment, structured force correction, deployable input constraints, and channel-specific output semantics. The comparison is intentionally log-conditioned: every model is evaluated against the reconstructed effective-wrench labels on held-out logs, not by replaying the vehicle trajectory in closed-loop simulation.

\subsection{Train/Validation/Test and Cross-Validation}
\label{subsec:split_protocol}

The primary experiment uses the fixed whole-log split described in Sec.~\ref{subsec:canonical_samples}. Prior alignment and correction fitting use training logs only. Model families and ridge regularization values are selected using validation logs, and the selected models are evaluated once on held-out test logs. Normalization statistics, affine prior-alignment parameters, feature fill values, and model-selection decisions are not fitted on test logs.

In the measured-mass-properties rerun, the fixed paper split is the primary evidence. The test split contains four complete logs from four flight dates, so per-log metrics are also reported to show how much the fixed-split result varies across held-out log composition. A full date-stratified cross-validation rerun with the measured-mass-properties branch is left as an optional robustness check before final submission.

\subsection{Replay Diagnostics}
\label{subsec:replay_diagnostics}

Replay is used only as a diagnostic for label construction, not as a primary validation metric. Outdoor open-loop replay compounds state-estimation noise, smoothing choices, time alignment, wind and disturbance terms, and controller-response mismatch. A corrected instantaneous effective-wrench map can be useful for simulation even when strict open-loop replay of the logged raw trajectory is not possible.

We therefore run alpha-only replay checks to verify the rotational label pathway. Moment-inferred angular acceleration is compared against the smoothed angular acceleration used for label construction, and the smoothed acceleration is integrated over short windows to quantify how closely it returns to the raw logged gyro. These checks diagnose the smoothing and label definitions; they do not validate closed-loop simulator behavior.

\subsection{Metrics}
\label{subsec:metrics}

The primary metrics are per-target root-mean-square error (RMSE) and mean absolute error (MAE):
\begin{equation}
  \mathrm{RMSE}_{j}
  =
  \sqrt{
    \frac{1}{N}
    \sum_{n=1}^{N}
    \left(
      \hat{y}_{n,j} - y_{n,j}
    \right)^2
  },
  \label{eq:rmse}
\end{equation}
\begin{equation}
  \mathrm{MAE}_{j}
  =
  \frac{1}{N}
  \sum_{n=1}^{N}
  \left|
    \hat{y}_{n,j} - y_{n,j}
  \right|.
  \label{eq:mae}
\end{equation}
Aggregate RMSE and MAE are reported as compact summaries, but per-target values are the main evidence because the six wrench channels have different units, amplitudes, excitation levels, and label quality.

To give the absolute errors a scale reference, we also report the held-out
standard deviation $\sigma_{y,j}$ of each target channel and the normalized
error $\mathrm{RMSE}_j/\sigma_{y,j}$. We report $R^2$ as a secondary metric.
Because the six channels have different scales and excitation levels, $R^2$ is
interpreted together with RMSE, MAE, and $\mathrm{RMSE}/\sigma_y$. This is
especially important for moment channels, where the label variance is much
smaller than in the force channels.
